\chapter{Harmonic Oscillator Review, Spin-1/2, and the First Exchange Principle}

These notes follow Leonard Susskind's fourth lecture in the supplementary \emph{Advanced Quantum Mechanics} sequence, with curation by LazyingArt LLC. The lecture has a very clear arc: a brief but purposeful return to the harmonic oscillator as preparation for field theory, a rapid reinterpretation of the two-state spin system as angular momentum, an atomic-physics motivation for spin and exclusion, and finally an intentionally unfinished first statement of identical-particle exchange. Throughout the oscillator discussion we follow the lecture convention \(m=1\) and, unless stated otherwise, \(\hbar=1\).

\section{Why the Harmonic Oscillator Returns}

Susskind does not begin with formalism. He begins by reminding the audience why the harmonic oscillator keeps coming back. Whenever a system has an equilibrium configuration and is displaced slightly from it, the motion is generically oscillatory. A violin string is a familiar example, but the real target is more fundamental: the electromagnetic field in a cavity, radiation in a waveguide, and lattice vibrations in a crystal all decompose into oscillatory modes.

This matters because the equilibrium configuration of the electromagnetic field is just the vacuum, namely no electric and magnetic field at all. Once that equilibrium is disturbed, the field oscillates. Likewise, a crystal supports collective vibrational modes, whose quanta are phonons. The harmonic oscillator is therefore not a detachable exercise from elementary quantum mechanics; it is the local mathematical model for a large part of quantum theory and, later, quantum field theory.

\section{Oscillator Algebra and the Quantized Spectrum}

The algebraic recap begins with the Hamiltonian
\begin{equation}
H=\frac{1}{2}p^{2}+\frac{1}{2}\omega^{2}x^{2}.
\end{equation}
Here \(\omega\) plays the role both of the oscillator frequency and of the parameter setting the curvature of the quadratic potential.

The ladder operators are defined by
\begin{equation}
a_{\pm}=\frac{p\pm i\omega x}{\sqrt{2\omega}}.
\end{equation}
Susskind emphasizes the normalization: the factor \(\sqrt{2\omega}\) is chosen precisely so that the commutator comes out in the simplest possible form,
\begin{equation}
[a_{-},a_{+}]=1.
\end{equation}
One then defines the number operator by
\begin{equation}
N=a_{+}a_{-}.
\end{equation}
The commutator implies that \(a_{+}\) and \(a_{-}\) raise and lower the eigenvalues of \(N\) by one unit, so the spectrum of \(N\) is quantized in nonnegative integers:
\[
0,1,2,3,\ldots
\]

If one ignores the additive constant, the Hamiltonian is simply
\begin{equation}
H=\omega N.
\end{equation}
Quantum mechanically, however, the noncommutativity of \(a_{+}\) and \(a_{-}\) produces the familiar correction
\begin{equation}
H=\omega\left(N+\frac{1}{2}\right).
\end{equation}
The extra \(\tfrac{1}{2}\omega\) is the ground-state energy, or in field-theory language the vacuum energy. Susskind explicitly remarks that for many purposes one can cross it out, because adding a constant to the energy does not change the equations of motion, but the constant is genuinely there.

At this point the lecture pauses over a basic question: why does the ladder begin at \(0\) at all? The answer is first given abstractly. Since \(H\) is built from \(p^{2}+\omega^{2}x^{2}\), it cannot have arbitrarily negative eigenvalues. There must therefore be a bottom state, and the only way a lowest state can avoid being lowered further is
\begin{equation}
a_{-}|0\rangle=0.
\end{equation}
Applying \(N=a_{+}a_{-}\) to this state gives
\begin{equation}
N|0\rangle=0.
\end{equation}
Thus the bottom state has number eigenvalue \(0\), even though its energy is not zero but \(\tfrac{1}{2}\omega\).

\begin{remark}
The label \(0\) on \(|0\rangle\) or \(\psi_{0}\) is only a name for the lowest state. It does not mean that the state vector itself is the zero vector.
\end{remark}

\section{Wave Functions, Square Integrability, and the Ground State}

A student question then triggers the main shift in representation. Susskind places the wave-function description alongside the ladder-operator description. The abstract state vector becomes a function of position, and energy eigenstates are written as \(\psi_{n}(x)\). The operator \(X\) becomes multiplication by \(x\), while the momentum becomes
\begin{equation}
X\mapsto x,\qquad P\mapsto -i\frac{d}{dx}.
\end{equation}
He remarks that \(\hbar\) can be restored later by dimensional analysis.

The time-independent Schr\"odinger equation is
\begin{equation}
\left[-\frac{1}{2}\frac{d^{2}}{dx^{2}}+\frac{1}{2}\omega^{2}x^{2}\right]\psi(x)=E\psi(x).
\end{equation}
This is a second-order differential equation, and for any \(E\) one can formally solve it. But most formal solutions are not acceptable physical states: they blow up at infinity and cannot be normalized. The relevant Hilbert space consists of square-integrable functions,
\begin{equation}
\int_{-\infty}^{\infty}\psi^{*}(x)\psi(x)\,dx=1.
\end{equation}
Only for special values of \(E\) does the differential equation admit such honest eigenfunctions. Susskind later notes that this restriction is also what makes \(a_{+}\) and \(a_{-}\) into Hermitian conjugates of one another.

\paragraph{Worked derivation of the ground state.}
Rather than solve the second-order equation directly, Susskind uses the algebra. The ground state must satisfy
\begin{equation}
a_{-}\psi_{0}(x)=0.
\end{equation}
Since overall numerical factors are irrelevant in a homogeneous equation, it is enough to use
\[
a_{-}\propto p-i\omega x.
\]
Substituting \(p=-i\,d/dx\) gives
\begin{equation}
(p-i\omega x)\psi_{0}(x)=0
\qquad\Longrightarrow\qquad
\left(\frac{d}{dx}+\omega x\right)\psi_{0}(x)=0.
\end{equation}
This is now a first-order equation. The standard trick is to write
\begin{equation}
\psi_{0}(x)=e^{f(x)}.
\end{equation}
Then \(\psi_{0}'=f'e^{f}\), and the common factor \(e^{f(x)}\) drops out, leaving
\begin{equation}
\frac{df}{dx}+\omega x=0.
\end{equation}
Integrating,
\begin{equation}
f(x)=-\frac{1}{2}\omega x^{2}+C,
\end{equation}
so the ground-state wavefunction is
\begin{equation}
\psi_{0}(x)\propto e^{-\omega x^{2}/2}.
\end{equation}
This is the Gaussian that one expects for the bottom of a quadratic potential. Its rapid decay makes it manifestly square integrable. The energy is already known from the algebra:
\begin{equation}
E_{0}=\frac{\omega}{2}.
\end{equation}
If one wants the normalized form, it is
\begin{equation}
\psi_{0}(x)=\left(\frac{\omega}{\pi}\right)^{1/4}e^{-\omega x^{2}/2}.
\end{equation}

\section{The First Excited State, Higher States, and the Classical Limit}

The first excited state is obtained in exactly the way the ladder algebra suggests:
\begin{equation}
\psi_{1}(x)\propto a_{+}\psi_{0}(x).
\end{equation}
Since \(a_{+}\propto p+i\omega x\), this becomes
\begin{equation}
\psi_{1}(x)\propto \left(-i\frac{d}{dx}+i\omega x\right)e^{-\omega x^{2}/2}.
\end{equation}
Using the ground-state relation already found, one sees immediately that the derivative term and the multiplication term reinforce rather than cancel. Up to an unimportant constant,
\begin{equation}
\psi_{1}(x)\propto x\,e^{-\omega x^{2}/2}.
\end{equation}

This form already contains the important physics. The first excited state is odd:
\[
\psi_{1}(-x)=-\psi_{1}(x).
\]
It vanishes at the origin, so the probability density \(|\psi_{1}(x)|^{2}\) also vanishes there. In other words, if the oscillator is in the first excited state, the probability of finding it exactly at the equilibrium point is zero.

Higher states are described only qualitatively in the lecture, but the qualitative structure is sharp. Each successive application of \(a_{+}\) produces a higher-degree polynomial multiplying the same Gaussian envelope. The parity alternates between even and odd. Each level adds a node. The wavefunction is pushed farther outward in \(x\), reflecting larger potential energy, while the increasingly rapid oscillations reflect larger momentum. Susskind's language here is vivid: more energy means both more support far from the origin and more ``wiggles'' near the center.

The lecture then turns to the correspondence principle. A single high-energy eigenstate is not itself a classical oscillator tracing out a trajectory. To see classical motion one must take a wave packet, that is, a superposition of many energy eigenstates. A displaced Gaussian packet is no longer an energy eigenstate; under the time-dependent Schr\"odinger equation its components acquire different phases, and the packet swings back and forth in a way that resembles the classical motion more and more closely at high energy. The classical limit emerges not from one stationary eigenfunction but from the time evolution of a suitable superposition.

\section{Spin-1/2 as Angular Momentum}

The lecture's second major subject begins with a change of interpretation rather than a new mathematical formalism. The two-state spin system discussed much earlier is now re-read as a genuine angular-momentum system. The key fact about angular momentum is not merely that it has three components, but that those components are associated with the spatial directions \(x\), \(y\), and \(z\), and that they satisfy the angular-momentum commutation relations
\begin{equation}
[L_{z},L_{x}]=iL_{y},
\end{equation}
together with cyclic permutations.

Susskind then compares these relations to the Pauli matrices,
\begin{equation}
\sigma_{z}=
\begin{pmatrix}
1&0\\
0&-1
\end{pmatrix},
\qquad
\sigma_{x}=
\begin{pmatrix}
0&1\\
1&0
\end{pmatrix},
\qquad
\sigma_{y}=
\begin{pmatrix}
0&-i\\
i&0
\end{pmatrix}.
\end{equation}
He chooses the \(z\)-representation, so \(\sigma_{z}\) is diagonal. Its eigenvalues are \(+1\) and \(-1\), corresponding to up and down.

The board calculation checks the commutator explicitly:
\begin{align}
\sigma_{z}\sigma_{x}
&=
\begin{pmatrix}
1&0\\
0&-1
\end{pmatrix}
\begin{pmatrix}
0&1\\
1&0
\end{pmatrix}
=
\begin{pmatrix}
0&1\\
-1&0
\end{pmatrix}
=
i\sigma_{y},\\
\sigma_{x}\sigma_{z}
&=
\begin{pmatrix}
0&1\\
1&0
\end{pmatrix}
\begin{pmatrix}
1&0\\
0&-1
\end{pmatrix}
=
\begin{pmatrix}
0&-1\\
1&0
\end{pmatrix}
=
-i\sigma_{y}.
\end{align}
Therefore
\begin{equation}
[\sigma_{z},\sigma_{x}]=2i\sigma_{y}.
\end{equation}
This is almost the angular-momentum algebra, but not quite: there is an extra factor of \(2\). The fix is simply a normalization,
\begin{equation}
S_{i}=\frac{\sigma_{i}}{2}.
\end{equation}
Then
\begin{equation}
[S_{z},S_{x}]=iS_{y},
\end{equation}
again with cyclic permutations. Thus the spin-\(\tfrac{1}{2}\) system is a perfectly good angular-momentum system.

In this normalization
\begin{equation}
S_{z}=
\frac{1}{2}
\begin{pmatrix}
1&0\\
0&-1
\end{pmatrix}
=
\begin{pmatrix}
\tfrac12&0\\
0&-\tfrac12
\end{pmatrix},
\end{equation}
so the eigenvalues of \(S_{z}\) are
\begin{equation}
m_{s}=\pm\frac{1}{2}.
\end{equation}
This is precisely the half-spin doublet. Susskind emphasizes that such spin can be regarded either as a literal tiny spinning object or as an abstract intrinsic property; whichever picture one prefers, its mathematical role is angular momentum.

\section{Atoms, Shell Counting, and the Need for Spin}

The next transition is from algebra to atomic structure. There are two kinds of angular momentum: orbital angular momentum and spin. For a particle one writes
\begin{equation}
\mathbf{J}=\mathbf{L}+\mathbf{S}.
\end{equation}
The orbital part is \(\mathbf{L}=\mathbf{R}\times\mathbf{P}\), while \(\mathbf{S}\) is intrinsic. Susskind concentrates on half-spin particles such as electrons, muons, neutrinos, and quarks.

The atomic motivation begins with hydrogen. In the lecture convention one writes
\begin{equation}
L^{2}=\ell(\ell+1),
\end{equation}
and for fixed orbital angular momentum \(\ell\) the multiplicity is
\begin{equation}
2\ell+1.
\end{equation}
The hydrogen spectrum is then drawn schematically as an \(E\)-versus-\(\ell\) diagram. At \(\ell=0\) there is a tower of radial excitations. The first \(\ell=1\) level is degenerate with the second \(\ell=0\) level; the first \(\ell=2\) level is degenerate with the third \(\ell=0\) level; and so on.

\begin{figure}
\begin{tikzpicture}[x=1.15cm,y=1.0cm]
\draw[->] (-0.4,0) -- (-0.4,5.2);
\draw[->] (-0.8,0) -- (5.8,0);
\node at (-0.75,5.0) {$E$};
\node at (5.6,-0.25) {$\ell$};

\node at (0,-0.35) {$0$};
\node at (2.5,-0.35) {$1$};
\node at (5.0,-0.35) {$2$};

\draw (-0.05,0.6) -- (0.35,0.6);
\draw (-0.05,1.8) -- (0.35,1.8);
\draw (-0.05,3.0) -- (0.35,3.0);
\draw (-0.05,3.9) -- (0.35,3.9);
\draw (-0.05,4.5) -- (0.35,4.5);

\draw (2.15,1.8) -- (2.35,1.8);
\draw (2.45,1.8) -- (2.65,1.8);
\draw (2.75,1.8) -- (2.95,1.8);

\draw (2.15,3.0) -- (2.35,3.0);
\draw (2.45,3.0) -- (2.65,3.0);
\draw (2.75,3.0) -- (2.95,3.0);

\draw (2.15,3.9) -- (2.35,3.9);
\draw (2.45,3.9) -- (2.65,3.9);
\draw (2.75,3.9) -- (2.95,3.9);

\draw (4.4,3.0) -- (4.55,3.0);
\draw (4.65,3.0) -- (4.8,3.0);
\draw (4.9,3.0) -- (5.05,3.0);
\draw (5.15,3.0) -- (5.3,3.0);
\draw (5.4,3.0) -- (5.55,3.0);

\draw (4.4,3.9) -- (4.55,3.9);
\draw (4.65,3.9) -- (4.8,3.9);
\draw (4.9,3.9) -- (5.05,3.9);
\draw (5.15,3.9) -- (5.3,3.9);
\draw (5.4,3.9) -- (5.55,3.9);

\node at (1.2,4.9) {$L^2=\ell(\ell+1)$};
\node at (1.0,4.45) {$2\ell+1$};
\end{tikzpicture}
\caption{Schematic hydrogenic degeneracy pattern used in the lecture. The vertical placement is not to scale; only the pattern of degeneracies and multiplicities matters here.}
\end{figure}

The resulting cumulative shell counts are
\begin{equation}
1,\qquad 1+3=4,\qquad 1+3+5=9,\qquad 1+3+5+7=16,
\end{equation}
that is,
\begin{equation}
1,4,9,16,\ldots=n^{2}.
\end{equation}
This is only the orbital counting. If each spatial state can be occupied by two opposite spin states, the idealized pattern doubles to
\begin{equation}
2,8,18,32,\ldots = 2n^{2}.
\end{equation}

This is where the periodic table enters. Helium is reasonably described, in a first approximation, by putting two electrons into the lowest spatial state. But lithium is not obtained by putting a third electron into that same state. The third electron goes into the first excited shell. Pauli's insight was that two electrons cannot occupy the same full quantum state. The reason helium still accommodates two electrons in its lowest orbital is that the electron has another two-valued property. Once one introduces spin-up and spin-down, the two electrons in the same spatial orbital are no longer in the same total state.

The shell counting in the lecture is deliberately naive, but it is very effective as motivation. The first excited shell contains one \(\ell=0\) state and three \(\ell=1\) states, for a total of four spatial states. Doubling by spin gives eight possibilities, which is exactly the number needed to understand the second row of the periodic table in this simplified picture.

Susskind also stresses that this shell model cannot be taken too literally. Once many-electron atoms are considered, electron-electron interactions become essential and the simple hydrogenic picture quickly ceases to be quantitatively reliable. Here the shell structure serves as a conceptual route to spin and exclusion, not as a finished many-body theory.

\section{Exchange, Identical Particles, and the First Statistics Principle}

The lecture's final move is to reframe exclusion as part of a more general question: what does it mean for quantum particles to be identical? Historically, one could treat exclusion as an extra postulate, but Susskind wants a cleaner formulation.

He first contrasts classical and quantum counting. In classical statistical mechanics one may imagine two identical particles carrying invisible labels; counting with or without those labels leads to a conceptual distinction, even if the thermodynamic answers come out the same after dividing by \(n!\). Quantum mechanics is sharper. Two electrons do not come with hidden labels. Interchanging them is an operation on the state, not a relabeling of two distinguishable objects.

For one particle the wavefunction is \(\psi(x)\). For two identical particles the natural description is
\begin{equation}
\psi(x_{1},x_{2}),
\end{equation}
where \(x_{1}\) and \(x_{2}\) refer to the two particle slots, not to two spatial directions. The probability density \(|\psi(x_{1},x_{2})|^{2}\) gives the probability of finding one particle at \(x_{1}\) and the other at \(x_{2}\).

To avoid confusion with the spin operator \(\mathbf{S}\), let us write the swap operator as \(P_{12}\), although Susskind simply called it \(S\). Its action is
\begin{equation}
P_{12}\psi(x_{1},x_{2})=\psi(x_{2},x_{1}).
\end{equation}
Applying the swap twice brings the state back to itself, so
\begin{equation}
P_{12}^{2}=1.
\end{equation}
Because it is an operation on the Hilbert space, \(P_{12}\) is unitary. Therefore its eigenvalues must satisfy \(\lambda^{2}=1\), and hence
\begin{equation}
\lambda=\pm 1.
\end{equation}
So identical-particle states must satisfy one of the two relations
\begin{equation}
P_{12}\psi=\psi,
\qquad\text{or}\qquad
P_{12}\psi=-\psi.
\end{equation}
This is exactly where the lecture stops. The wavefunction of a pair of identical particles either comes back to itself under exchange or comes back with a minus sign. The full development of bosons and fermions is deferred to the next lecture; here the classification principle is posed, but not yet unpacked.

\section{Summary}

The lecture begins with the ubiquity of the harmonic oscillator and uses that fact to justify a rapid return to its algebra. From the ladder operators one obtains the quantized spectrum, the ground-state condition, and then, in position space, the Gaussian ground state and the odd first excited state. The discussion then broadens into the semiclassical meaning of higher states and wave packets. From there the two-state spin system is reinterpreted as a spin-\(\tfrac12\) angular momentum system through the Pauli matrices and their commutators. Atomic shell counting then supplies the physical motivation for spin and for Pauli's exclusion idea. Finally, exclusion is recast as the first step toward a theory of identical particles: the state of two identical particles is classified by what happens when the particles are exchanged.